\newif\ifglobal

%%%%%%%%%%%%%%%%%%%%%%%
%%% Compile to view %%%
%%%%%%%%%%%%%%%%%%%%%%%

%%%%%%%%%%%%%%%%%%%%%%%%%%%%%%%%%%%%%%%%%%%%%%%%%%%
%%% Readme:                                     %%%
%%% https://www.overleaf.com/read/bwdjvfjksvjy/ %%%
%%%%%%%%%%%%%%%%%%%%%%%%%%%%%%%%%%%%%%%%%%%%%%%%%%%

% >>> M?: marks a module setting number or important notice

% >>> M4: Set {./relative-path} to external files for this module <<<
\def \modulefiles {./15-WDT}
% To add an external file, use:
% (Replace '---' with actual filename)
% '\input{\modulefiles/tables/---.tex}' for tables
% '\includegraphics{\modulefiles/figures/---}' for figures
% '\subfile{\modulefiles/---}' for register files

% >>> M1: This module is in global project? <<<
% 'YES' - uncomment; 'NO' - comment out
\globaltrue

\ifglobal
    \documentclass[main\_\_CN.tex]{subfiles}

\else
    % Font "FandolSong-Regular" used by ctex does not have GSUB table.
    % It causes warnings that do not affect compilation and can be ignored.
    \PassOptionsToPackage{quiet}{fontspec}

    \documentclass[openany, 10pt]{book}

    % >>> M2: In ./00-shared/preamble-trm-module, also find
    %         settings for visibility of todo notes, watermark <<<
    \usepackage{preamble-trm-module}


    % Enable Chinese typesetting
    \usepackage{ctex}

    % >>> M3: Temporary labels in Chinese
    %         you can add your own labels there <<<
    \myexternaldocument{./00-shared/temp-labels-cn}

\fi %\ifglobal


\setcounter{secnumdepth}{4}

% >>> M5: If this module needs any updates in
% './00-shared/preamble-trm-module.sty'
% (include additional package, etc.), contact your module PM
% or create the file './00-shared/changelog.tex' make the
% necessary updates yourself and document them in this file <<<


\begin{document}


% Sets language into which pre-defined bits of text are translated
\selectlanguage{chinese}

% Do not delete this block
\ifglobal
% Add the following front matter if
% this module is outside of global project
\else
    \InputIfFileExists{readme}{}{}
    {\let\clearpage\relax \listoftodos}
    \tableofcontents
    %\listoftables
    %\listoffigures
\fi


%%%%%%%%%%%%%%%%%%%%%%%%%
%%% TRM Module Begins %%%
%%%%%%%%%%%%%%%%%%%%%%%%%

\hypertarget{wdt}{}
\chapter{看门狗定时器 (WDT)}
\label{mod:wdt}

\section{概述}

看门狗定时器是一种硬件定时器，用于检测和修复故障。软件必须定期喂狗（复位），以防超时。系统或软件若出现不可预知的问题（比如软件卡在某个循环或逾期事件中）将无法按时喂狗，造成看门狗超时。因此，看门狗定时器有助于检测、处理系统或软件的错误行为。

如图 \ref{fig:wdt-timers-overview} 所示，\chipname{} 中有三个数字看门狗定时器：章节 \ref{mod:timg} \textit{\nameref{mod:timg}} 描述的两个定时器组中各有一个（称作主系统看门狗定时器，缩写为 MWDT），RTC 模块中有一个（称作 RTC 看门狗定时器，缩写为 RWDT）。数字看门狗在运行期间会经历四个阶段（除非看门狗按时喂狗或者处于关闭状态），每个阶段均可配置单独的超时时间和超时动作，其中 MWDT 支持中断、CPU 复位和内核复位三种超时动作，RWDT 支持中断、CPU 复位、内核复位和系统复位四种超时动作（详见章节 \ref{sec:wdt-timeout} \textit{\nameref{sec:wdt-timeout}}）。每个阶段的超时时间都可单独设置。

在 flash 引导模式下，RWDT 和定时器组 0 的 MWDT 会默认使能，以检测引导过程中发生的错误，并恢复运行。

\chipname{} 中还有一个模拟看门狗定时器——超级看门狗 (SWD)。超级看门狗是模拟域的超低功耗电路，可以防止系统在数字电路异常状态下运行，并在必要时复位系统。

\begin{figure}[H]
    \centering
    \includegraphics[width=0.5\textwidth]{\modulefiles/figures/esp32-s3-wdt-overview}
    \caption{看门狗定时器概览}
    \label{fig:wdt-timers-overview}
\end{figure}

请注意，本章节仅包含看门狗定时器的功能描述，其寄存器部分详见章节 \ref{mod:timg} \textit{\nameref{mod:timg}} 和章节 \ref{mod:lowpowm} \textit{\nameref{mod:lowpowm}}。

\section{数字看门狗定时器}
\subsection{主要特性}
看门狗定时器具有如下特性：
\begin{itemize}
    \item 四个阶段，每个阶段都可配置超时时间。每阶段都可单独配置、使能和关闭
    \item 如在某个阶段发生超时，MWDT 会采取中断、CPU 复位和内核复位中的一种超时动作，RWDT 则会采取中断、CPU 复位、内核复位和系统复位中的一种超时动作
    \item 32 位超时计数器
    \item 写保护，防止 RWDT 和 MWDT 配置误改动
    \item flash 启动保护\\
    如果在预定时间内 SPI flash 的引导过程没有完成，看门狗会重启整个主系统
\end{itemize}

\subsection{功能描述}

\begin{figure}[H]
    \centering
    \includegraphics[width=0.9\textwidth]{\modulefiles/figures/725wdt_20210115}
    \caption{\chipname{} 的看门狗定时器}
    \label{fig:wdt-timers}
\end{figure}

图 \ref{fig:wdt-timers} 为 \chipname{} 数字系统中的三个看门狗定时器。

\subsubsection{时钟源与 32 位计数器}

每个看门狗定时器的核心是一个 32 位计数器。APB 时钟经过可配置的 16 位预分频器后会得到 MWDT 的时钟源。RWDT 的时钟源则直接取自于 RTC 慢速时钟（RTC 慢速时钟源详见章节 \ref{mod:resclk} \textit{\nameref{mod:resclk}}）。MWDT 的 16 位预分频器可通过 \hyperref[regdesc:TIMGWDTCONFIG1REG]{TIMG\_WDTCONFIG1\_REG} 寄存器的 \hyperref[fielddesc:TIMGWDTCLKPRESCALE]{TIMG\_WDT\_CLK\_PRESCALE} 字段配置。

MWDT 和 RWDT 看门狗可分别通过设置 \hyperref[fielddesc:TIMGWDTEN]{TIMG\_WDT\_EN} 和 \hyperref[fielddesc:RTCCNTLWDTEN]{RTC\_CNTL\_WDT\_EN} 字段使能。看门狗使能后，其内部 32 位计数器的值会在每个时钟源周期内累加 1，直到达到该阶段的超时时间（即在该阶段发生超时）。如发生超时，计数器的值会重置为 0，同时看门狗进入下一阶段。如果软件在规定的时间内成功喂狗，看门狗定时器会回到阶段 0，并将计数器的值重置为 0。软件向 \hyperref[regdesc:TIMGWDTFEEDREG]{TIMG\_WDTFEED\_REG} 和 \hyperref[fielddesc:RTCCNTLRTCWDTFEED]{RTC\_CNTL\_RTC\_WDT\_FEED} 寄存器内写入任意值，便可分别为 MDWT 和 RWDT 喂狗。


\subsubsection{阶段与超时动作}\label{sec:wdt-timeout}
定时器在各阶段可以配置不同的超时时间和对应的超时动作。某一阶段超时会触发对应的超时动作，同时计数器的值被重置为 0，看门狗进入下一阶段。MWDT 和 RWDT 有四个阶段（称为阶段 0 至阶段 3）。看门狗定时器会循环工作（即从阶段 0 至阶段 3，再回到阶段 0）。

MWDT 每个阶段的超时时间可用 \hyperref[regdesc:TIMGWDTCONFIG2REG]{TIMG\_WDTCONFIG\textcolor{red}{\it{i}}\_REG}（i 的范围是 2 到 5）寄存器配置，RWDT 的超时时间可用 \hyperref[fielddesc:RTCCNTLWDTSTG3]{RTC\_CNTL\_WDT\_STG\textcolor{red}{\it{j}}\_HOLD}（j 的范围是 0 到 3）字段配置。

值得注意的是，RWDT 在阶段 0 的超时时间 (T\textsubscript{hold0}) 受 eFuse 寄存器 \hyperref[regdesc:EFUSERDREPEATDATA1REG]{EFUSE\_RD\_REPEAT\_DATA1\_REG} 的 \\\hyperref[fielddesc:EFUSEWDTDELAYSEL]{EFUSE\_WDT\_DELAY\_SEL} 字段和 \hyperref[fielddesc:RTCCNTLWDTSTG0HOLD]{RTC\_CNTL\_WDT\_STG0\_HOLD} 字段共同影响，关系如下：

\begin{displaymath}
T_{hold0} = \hyperref[fielddesc:RTCCNTLWDTSTG0HOLD]{RTC\_CNTL\_WDT\_STG0\_HOLD}   << (\hyperref[fielddesc:EFUSEWDTDELAYSEL]{EFUSE\_WDT\_DELAY\_SEL} + 1)
\end{displaymath}

其中，$<<$ 为左移运算符。

如某个阶段超时，下列超时动作之一将会执行：
\begin{itemize}
    \item 触发中断\\
    如阶段超时，中断被触发。
    \item CPU 复位 -- 复位 CPU 核心\\
    如阶段超时，复位 CPU 核心。
    \item 内核复位 -- 复位主系统\\
    如阶段超时，主系统（包括 MWDT、CPU 和所有外设）复位。功耗管理单元和 RTC 外设不会复位。
    \item 系统复位 -- 复位主系统、功耗管理单元和 RTC 外设\\
    如阶段超时，主系统、功耗管理单元和 RTC 外设（详见章节 \ref{mod:lowpowm} \textit{\nameref{mod:lowpowm}}）同时复位。此动作仅可在 RWDT 中实现。
    \item 关闭\\
    该阶段对系统不产生影响。
\end{itemize}

MWDT 所有阶段的超时动作均在 \hyperref[regdesc:TIMGWDTCONFIG0REG]{TIMG\_WDTCONFIG0\_REG} 寄存器中配置。RWDT 的超时动作可在 \\\hyperref[regdesc:RTCCNTLRTCWDTCONFIG0REG]{RTC\_CNTL\_WDTCONFIG0\_REG} 寄存器配置。

\subsubsection{写保护}
看门狗定时器对于检测和处理系统或软件错误而言至关重要，不应轻易关闭（例如，因写寄存器位置错误而误将看门狗关闭）。因此，MWDT 和 RWDT 引入写保护机制，防止看门狗因无意的写操作而被关闭或篡改。

写保护机制通过每个看门狗定时器的写密钥字段运行（MWDT 看门狗使用 \hyperref[fielddesc:TIMGWDTWKEY]{TIMG\_WDT\_WKEY}，RWDT 看门狗使用 \hyperref[fielddesc:RTCCNTLWDTWKEY]{RTC\_CNTL\_WDT\_WKEY}）。必须向看门狗定时器的写密钥字段写入 0x{}50D83AA1，才能修改其它看门狗寄存器。如果写密钥字段的值不是 0x{}50D83AA1，任何试图向看门狗定时器寄存器（除了向写密钥字段本身）写值的操作都会被忽略。推荐按以下步骤访问看门狗定时器：

\begin{enumerate}
\item 将 0x{}50D83AA1 写入看门狗定时器的写密钥字段，关闭写保护。
\item 根据需要修改看门狗，如喂狗或改变配置。
\item 向看门狗定时器的写密钥字段上写入除 0x{}50D83AA1 以外的任意值，重新使能写保护。
\end{enumerate}

\subsubsection{flash 引导保护}
在 flash 引导模式下，定时器组 0（见图 \ref{fig:timg} \textit{\nameref{fig:timg}}）的 MWDT 和 RWDT 会默认使能。MWDT 的阶段 0 的默认超时动作为内核复位（复位主系统）。RWDT 的阶段 0 超时动作为系统复位（复位主系统和 RTC）。引导后，应将 \hyperref[fielddesc:TIMGWDTFLASHBOOTMODEN]{TIMG\_WDT\_FLASHBOOT\_MOD\_EN} 和 \hyperref[fielddesc:RTCCNTLWDTFLASHBOOTMODEN]{RTC\_CNTL\_WDT\_FLASHBOOT\_MOD\_EN} 位清零，分别关闭 MWDT 和 RWDT 的 flash 引导保护。然后，软件可以配置 MWDT 和 RWDT。

\section{模拟看门狗定时器}\label{sec:wdt-swdt}
超级看门狗 (SWD) 是模拟域的超低功耗电路，可以防止系统在数字电路异常状态下运行，并在必要时复位系统。SWD 包含一个看门狗电路，需在每个超时阶段（约不足一秒）至少喂狗一次。该电路会在看门狗超时时间约 100 ms 之前发送 WD\_INTR 信号提醒系统喂狗。

如果系统不回应 SWD 的喂狗请求，看门狗超时，SWD 会产生系统电平信号 SWD\_RSTB，复位芯片上的整个数字电路。

\subsection{主要特性}
SWD 具有如下特性：
\begin{itemize}
    \item 超低功耗
    \item 用中断提醒 SWD 即将超时
    \item 软件有多种专用的方法喂 SWD，让 SWD 监控整个操作系统的工作状态
\end{itemize}

\subsection{SWD 控制器}
\subsubsection{结构}

\begin{figure}[H]
    \centering
    \includegraphics[width=0.9\textwidth]{\modulefiles/figures/swd_20210115}
    \caption{SWD 控制器结构}
    \label{fig:wdt}
\end{figure}

\subsubsection{工作流程}
正常状态下：
\begin{itemize}
    \item SWD 控制器收到 SWD 的喂狗请求。
    \item SWD 控制器可以向主 CPU 或 ULP-RISC-V 发送中断。
    \item 主 CPU 可以决定是通过置位 \hyperref[fielddesc:RTCCNTLSWDFEED]{RTC\_CNTL\_SWD\_FEED} 直接喂狗，还是发送中断让 ULP-RISC-V 置位该字段喂狗。
    \item CPU 或 ULP-RISC-V 喂狗时，需要先向 \hyperref[fielddesc:RTCCNTLSWDWKEY]{RTC\_CNTL\_SWD\_WKEY} 写 0x{}8F1D312A  关闭 SWD 控制器的写保护。这样做可以防止系统在数字电路异常状态下运行时误喂 SWD。
    \item 如将 \hyperref[fielddesc:RTCCNTLSWDAUTOFEEDEN]{RTC\_CNTL\_SWD\_AUTO\_FEED\_EN} 置 1，SWD 控制器也可配置为在不需要 CPU 或 ULP-RISC-V 干预的情况下喂 SWD。
\end{itemize}

复位后：
\begin{itemize}
    \item 可查看 \hyperref[fielddesc:RTCCNTLRESETCAUSEPROCPU]{RTC\_CNTL\_RESET\_CAUSE\_PROCPU}[5:0] 获知 CPU 复位原因。\\
	如 \hyperref[fielddesc:RTCCNTLRESETCAUSEPROCPU]{RTC\_CNTL\_RESET\_CAUSE\_PROCPU}[5:0] == 0x{}12，则表示上一次复位的原因是 SWD 复位。
	\item 置位 \hyperref[fielddesc:RTCCNTLSWDRSTFLAGCLR]{RTC\_CNTL\_SWD\_RST\_FLAG\_CLR} 清除 SWD 复位标志。
\end{itemize}

\section{中断}

看门狗定时器中断，请前往章节 \ref{mod:timg} \textit{\nameref{mod:timg}} 的第 \ref{sec:timg-interrupt} 节 \textit{\nameref{sec:timg-interrupt}} 查看。

\section{寄存器}
MWDT 寄存器是定时器组模块的一部分，在章节 \ref{mod:timg} \textit{\nameref{mod:timg}} 的第 \ref{sec:timg-reg-summ} 节 \textit{\nameref{sec:timg-reg-summ}} 中有详细描述。RWDT 和 SWD 寄存器是 RTC 模块的一部分，在章节 \ref{mod:lowpowm} \textit{\nameref{mod:lowpowm}} 的第 \ref{sec:lowpowm-reg-summ} 节 \textit{\nameref{sec:lowpowm-reg-summ}} 中有详细描述。


\end{document}
